\documentclass[a4paper,11pt]{article}
\usepackage[utf8]{inputenc}
\usepackage[T1]{fontenc}
\usepackage[french]{babel}
\usepackage{geometry}
\geometry{margin=2.2cm}
\usepackage{amsmath}
\usepackage{amssymb}
\usepackage{booktabs}
\usepackage{hyperref}
\usepackage{listings}
\usepackage{xcolor}
\usepackage{enumitem}

\definecolor{codebg}{HTML}{F5F7FA}
\definecolor{codeborder}{HTML}{C5D0DC}
\definecolor{codecomment}{HTML}{2E8B57}
\definecolor{codekeyword}{HTML}{0000FF}
\definecolor{codestring}{HTML}{A31515}

\lstdefinestyle{cstyle}{
  language=C++,
  backgroundcolor=\color{codebg},
  basicstyle=\ttfamily\small,
  keywordstyle=\color{codekeyword}\bfseries,
  commentstyle=\color{codecomment}\itshape,
  stringstyle=\color{codestring},
  frame=single,
  rulecolor=\color{codeborder},
  breaklines=true,
  showstringspaces=false
}

\title{Analyse FFT et traitement du signal\\dans \texttt{pc\_analyzer}}
\author{Projet Radar Doppler 24.125 GHz}
\date{Mars 2026}

\begin{document}
\maketitle
\tableofcontents
\newpage

\section{Objectif de ce document}

Ce document explique, en profondeur, la cha\^ine de traitement du signal du programme \texttt{pc\_analyzer} (\texttt{radar\_analyzer.cpp}), depuis l'acquisition ADC c\^ot\'e STM32 jusqu'au r\'esultat final affich\'e \`a l'op\'erateur (fr\'equence Doppler, vitesse, SNR, spectre, spectrogramme, tendance).

L'objectif est double :
\begin{itemize}[leftmargin=2em]
  \item \textbf{Technique} : comprendre exactement ce que fait le code.
  \item \textbf{Th\'eorique} : relier chaque choix d'impl\'ementation aux principes DSP/radar.
\end{itemize}

\section{Vue syst\`eme: de l'onde radar \`a la vitesse affich\'ee}

\subsection{Cha\^ine fonctionnelle}

\[
\text{Cible en mouvement}
\rightarrow
\text{D\'ecalage Doppler } f_d
\rightarrow
\text{Signal IF analogique}
\rightarrow
\text{ADC STM32 (100 kHz)}
\rightarrow
\text{Hops UDP}
\rightarrow
\text{Fen\^etre glissante PC}
\rightarrow
\text{FFT + d\'etection}
\rightarrow
\text{Estimation vitesse}
\rightarrow
\text{Affichage temps r\'eel}
\]

\subsection{Relation Doppler utilis\'ee}

Le code utilise le facteur :
\[
f_d \approx 161 \cdot v \quad (\text{Hz pour } v \text{ en m/s})
\]
soit :
\[
v = \frac{f_d}{161}
\]

Ce facteur provient de la formule monostatique :
\[
f_d = \frac{2 v f_0}{c}
\]
avec \(f_0 = 24.125\,\text{GHz}\).

\section{Param\`etres fondamentaux de \texttt{pc\_analyzer}}

Le code fixe les constantes de base suivantes :
\begin{lstlisting}[style=cstyle]
FS = 100000.0f
STM32_FRAME_SIZE = 4096
STM32_HOP_SIZE = 1024
UDP_PORT = 5555
DOPPLER_SCALE = 161.0f
\end{lstlisting}

Interpr\'etation physique :
\begin{itemize}[leftmargin=2em]
  \item Un hop de 1024 \'echantillons couvre \(1024/100000 = 10.24\) ms.
  \item Une fen\^etre de 4096 couvre \(40.96\) ms.
  \item Le moteur FFT PC est configurable (\(1024 \rightarrow 65536\) points), avec valeur par d\'efaut 8192.
\end{itemize}

\section{Transport UDP et reconstruction du signal}

\subsection{Fragmentation c\^ot\'e r\'eseau}

Chaque hop ADC (1024 \(\times\) uint16 = 2048 octets) est fragment\'e. C\^ot\'e PC, le thread UDP d\'ecode un en-t\^ete de 12 octets :
\begin{itemize}[leftmargin=2em]
  \item \texttt{seq} : num\'ero de s\'equence hop,
  \item \texttt{frag\_id}, \texttt{frag\_count},
  \item \texttt{offset}, \texttt{payload\_len}.
\end{itemize}

Le r\'eassemblage utilise un bitmap de fragments pour :
\begin{itemize}[leftmargin=2em]
  \item accepter les arriv\'ees hors ordre,
  \item ignorer les doublons,
  \item d\'etecter les hops incomplets/perdus.
\end{itemize}

\subsection{Fen\^etre glissante locale}

Une fois un hop complet re\c{c}u :
\begin{enumerate}[leftmargin=2em]
  \item on d\'ecale \texttt{rolling\_adc} de 1024 \'echantillons (\texttt{memmove}),
  \item on ins\`ere le nouveau hop en fin de fen\^etre.
\end{enumerate}

Cette op\'eration cr\'ee une analyse \textbf{overlap 75\%} :
\[
\text{pas temporel} = 1024, \quad \text{taille fen\^etre} = 4096
\]
ce qui am\'eliore la r\'eactivit\'e d'affichage sans sacrifier la stabilit\'e fr\'equentielle.

\section{Pr\'e-traitement temporel avant FFT}

\subsection{Notations et cadre math\'ematique}

Dans cette section, nous notons :
\begin{itemize}[leftmargin=2em]
  \item \(x[n]\) : \'echantillon ADC brut (apr\`es reconstruction de fen\^etre glissante),
  \item \(N_w = 4096\) : longueur physique de la fen\^etre glissante trait\'ee,
  \item \(N_{FFT}\) : taille de FFT demand\'ee (\(N_{FFT}\ge N_w\), puissance de 2),
  \item \(F_s = 100\,000\) Hz : fr\'equence d'\'echantillonnage.
\end{itemize}

\subsection{Suppression de la composante continue (DC removal)}

Le code estime la moyenne temporelle puis centre le signal :
\[
\mu = \frac{1}{N_w}\sum_{n=0}^{N_w-1}x[n],
\qquad
x_c[n]=x[n]-\mu.
\]

\paragraph{Pourquoi c'est crucial.}
Sans cette \'etape, la composante quasi-statique de l'ADC (offset analogique + offset num\'erique) concentre une \'energie importante autour du bin \(k=0\). Via fuite spectrale, cette \'energie se propage dans les bins voisins et r\'eduit fortement la sensibilit\'e de d\'etection aux faibles Doppler.

\paragraph{Lecture physique.}
On ne cherche pas l'\'energie absolue du front-end, mais la \emph{modulation Doppler}. Le centrage transforme donc le probl\`eme en ``analyser des variations autour d'un point d'\'equilibre''.

\subsection{Fen\^etre de Hann : forme, r\^ole, et compromis}

Le fen\^etrage est :
\[
w[n] = \frac{1}{2}\left(1-\cos\left(\frac{2\pi n}{N_w-1}\right)\right), \quad 0\le n < N_w
\]
et
\[
x_w[n]=x_c[n]\,w[n].
\]

\paragraph{Composants de la formule.}
\begin{itemize}[leftmargin=2em]
  \item Le terme \(\cos(\cdot)\) impose une transition douce aux extr\'emit\'es.
  \item Le facteur \(1/2\) normalise la plage vers \([0,1]\).
  \item La division par \(N_w-1\) aligne exactement les bords \(n=0\) et \(n=N_w-1\) \`a z\'ero.
\end{itemize}

\paragraph{Pourquoi fen\^etrer en radar Doppler.}
La FFT suppose implicitement un signal p\'eriodique sur la fen\^etre analys\'ee. Sans fen\^etre, la couture entre fin et d\'ebut cr\'ee une discontinuit\'e, ce qui \'etale l'\'energie d'une raie dans de nombreux bins (\emph{spectral leakage}). La Hann r\'eduit drastiquement les lobes secondaires, au prix d'un lobe principal plus large.

\paragraph{Point d'impl\'ementation important.}
Le code calcule \(w[n]\) sur \(N_w=4096\), m\^eme si \(N_{FFT}>4096\). C'est correct : la fen\^etre doit \'epouser la portion effectivement observ\'ee du signal, puis le padding intervient \emph{apr\`es}.

\subsection{Zero-padding : ce que c'est, et ce que ce n'est pas}

Si \(N_{FFT}>N_w\), le code applique :
\[
x_p[n]=
\begin{cases}
x_w[n], & 0\le n < N_w\\
0, & N_w \le n < N_{FFT}.
\end{cases}
\]

\paragraph{Interpr\'etation rigoureuse.}
Le zero-padding n'augmente pas l'information (pas de nouvelle mesure), mais augmente la densit\'e d'\'echantillonnage de la transform\'ee discr\`ete en fr\'equence. En pratique, cela facilite la lecture du pic, l'affichage, et la stabilit\'e des m\'ethodes de d\'etection bas\'ees sur le ``maximum de bin''.

\subsection{Exemple num\'erique (section 5)}

Prenons un bloc de 4096 \'echantillons ADC avec :
\[
\mu = 2042.7 \text{ LSB},\quad x[1000]=2190 \text{ LSB}.
\]
Apr\`es centrage :
\[
x_c[1000]=2190-2042.7=147.3 \text{ LSB}.
\]

Supposons \(w[1000]=0.31\) (Hann). Alors :
\[
x_w[1000]=147.3\times 0.31 \approx 45.7 \text{ LSB}.
\]

Si \(N_{FFT}=8192\), les 4096 \'echantillons suivants sont z\'ero :
\[
x_p[4096]=x_p[4097]=\cdots=x_p[8191]=0.
\]
On a donc la m\^eme information physique que sur 4096 points, mais un maillage spectral plus fin.

\section{FFT, magnitude et \'echelle en dB}

\subsection{Transform\'ee discr\`ete en fr\'equence}

Le moteur \texttt{kissfft::rfft} calcule :
\[
X[k]=\sum_{n=0}^{N_{FFT}-1}x_p[n]e^{-j2\pi kn/N_{FFT}},\quad 0\le k < N_{FFT}.
\]
Comme le signal est r\'eel, le spectre est conjugu\'e-sym\'etrique, et l'on exploite la moiti\'e positive (\(0\) \`a \(N_{FFT}/2\)).

\subsection{Correspondance bin \(\leftrightarrow\) fr\'equence}

Chaque bin \(k\) repr\'esente :
\[
f_k = k\frac{F_s}{N_{FFT}},\qquad \Delta f = \frac{F_s}{N_{FFT}}.
\]

\paragraph{Exemple.}
Pour \(F_s=100\,\text{kHz}\), \(N_{FFT}=8192\) :
\[
\Delta f \approx 12.207\,\text{Hz/bin}.
\]

\subsection{Module complexe et passage en dB}

Le code calcule :
\[
|X[k]| = \sqrt{\Re(X[k])^2+\Im(X[k])^2},
\]
puis
\[
S[k]_{\mathrm{dB}} = 20\log_{10}\!\big(|X[k]|+\varepsilon\big),\quad \varepsilon=10^{-12}.
\]

\paragraph{Pourquoi \(20\log_{10}\) et pas \(10\log_{10}\).}
\(S[k]\) est un \emph{module d'amplitude}. Pour une grandeur de puissance, on utiliserait \(10\log_{10}\). Ici on convertit une amplitude, d'o\`u le facteur 20.

\paragraph{R\^ole de \(\varepsilon\).}
Il garantit la d\'efinition num\'erique du logarithme m\^eme si \(|X[k]|\approx 0\), sans alt\'erer la dynamique utile.

\subsection{Lissage exponentiel bin par bin}

Le spectre liss\'e suit :
\[
S_s^{(m)}[k]=(1-\alpha)S_s^{(m-1)}[k]+\alpha S^{(m)}[k],
\]
o\`u \(m\) est l'indice temporel de hop trait\'e.

\paragraph{Constante de temps effective.}
Avec un pas temporel \(T_h=10.24\) ms (hop de 1024), une approximation de la constante de temps est :
\[
\tau \approx \frac{T_h}{\alpha}.
\]
Donc \(\alpha\) petit donne un spectre tr\`es stable (mais moins r\'eactif).

\subsection{Exemple num\'erique (section 6)}

Consid\'erons \(F_s=100000\) Hz.

\paragraph{Cas A : \(N_{FFT}=8192\).}
\[
\Delta f = \frac{100000}{8192}=12.207\ \text{Hz/bin}.
\]
Un pic au bin \(k=410\) correspond \`a :
\[
f_{410}=410\times 12.207 \approx 5005\ \text{Hz}.
\]

\paragraph{Conversion dB.}
Si \(|X[410]|=2500\), alors :
\[
S[410]_{dB}=20\log_{10}(2500)\approx 67.96\ \text{dB}.
\]
Si le frame pr\'ec\'edent avait \(S_s^{(m-1)}[410]=64\) dB et \(\alpha=0.25\) :
\[
S_s^{(m)}[410]=0.75\times64+0.25\times67.96
=64.99\ \text{dB}.
\]
On voit la stabilisation : le saut brut de \(+3.96\) dB devient \(+0.99\) dB.

\section{D\'etection CFAR et robustification bruit}

\subsection{CA-CFAR : d\'etail de la construction du seuil}

Pour une cellule test \(k\) (\emph{CUT}), on d\'efinit :
\begin{itemize}[leftmargin=2em]
  \item \(G\) cellules de garde de chaque c\^ot\'e (pour exclure l'\'energie du pic),
  \item \(R\) cellules de r\'ef\'erence de chaque c\^ot\'e (estimation bruit local).
\end{itemize}

Le nombre de cellules de r\'ef\'erence effectives vaut \(N_{ref}\approx 2R\) (hors bords).

Le bruit local estim\'e en dB :
\[
\widehat{N}[k]=\frac{1}{N_{ref}}
\left(\sum_{i\in\mathcal{L}_k}S_s[i]+\sum_{i\in\mathcal{R}_k}S_s[i]\right),
\]
avec \(\mathcal{L}_k,\mathcal{R}_k\) les ensembles de r\'ef\'erence gauche/droite.

Le seuil :
\[
T[k]=\widehat{N}[k]+\alpha_{CFAR},
\]
et la d\'etection binaire :
\[
D[k]=
\begin{cases}
1,& S_s[k]>T[k]\\
0,& \text{sinon}.
\end{cases}
\]

\paragraph{Sens de \(\alpha_{CFAR}\).}
\(\alpha_{CFAR}\) est une marge en dB au-dessus du bruit local. Plus \(\alpha\) augmente, plus la probabilit\'e de fausse alarme baisse, mais plus la probabilit\'e de d\'etection des cibles faibles diminue.

\subsection{Traitement des perturbations secteur et harmoniques}

Le code inclut un masque ``mains'' : bins proches de \(h f_{mains}\) (\(h=1..H\)) sont p\'enalis\'es/invalid\'es.

Math\'ematiquement, pour un bin \(k\), si
\[
|f_k-hf_{mains}|\le W_{mains},
\]
alors \(D[k]\leftarrow 0\) et \(T[k]\) est major\'e d'une marge suppl\'ementaire.

Interpr\'etation : on injecte une connaissance \emph{a priori} sur une source d'interf\'erence quasi-d\'eterministe.

\subsection{Notch param\'etrique}

Si activ\'e, un notch logiciel att\'enue les bins de la bande :
\[
\mathcal{B}_{notch}=\left[f_c-\frac{B}{2},\;f_c+\frac{B}{2}\right]
\]
avec profondeur \(D_{notch}\) (en dB). Concr\`etement, le code soustrait \(D_{notch}\) \`a \(S[k]\) et \(S_s[k]\) dans cette bande.

\subsection{Plancher de bruit adaptatif : mod\`ele dynamique}

Le bruit instantan\'e \(N_{inst}\) est la m\'ediane des bins non d\'etect\'es et non notched. Puis :
\[
N_f^{(m+1)}=
\begin{cases}
N_f^{(m)}+\alpha_{\uparrow}(N_{inst}^{(m)}-N_f^{(m)}), & N_{inst}^{(m)}>N_f^{(m)}\\
N_f^{(m)}+\alpha_{\downarrow}(N_{inst}^{(m)}-N_f^{(m)}), & N_{inst}^{(m)}\le N_f^{(m)}.
\end{cases}
\]

Ensuite on force :
\[
T[k]\ge N_f+\Delta_{offset}.
\]

\paragraph{Lecture ``syst\`eme''.}
\(\alpha_{\uparrow}\) rapide permet de suivre une hausse soudaine du bruit; \(\alpha_{\downarrow}\) plus lent \'evite de ``d\'ecrocher'' trop vite apr\`es un transitoire. C'est un m\'ecanisme d'hyst\'er\'esis temporelle sur le plancher.

\subsection{Exemple num\'erique (section 7)}

Prenons :
\[
G=4,\quad R=16,\quad \alpha_{CFAR}=12\ \text{dB}.
\]
Donc \(N_{ref}\approx 32\) cellules.

Pour un CUT \(k\), supposons :
\[
\widehat{N}[k]=-58.5\ \text{dB}.
\]
Le seuil devient :
\[
T[k]=-58.5+12=-46.5\ \text{dB}.
\]

Si \(S_s[k]=-43.0\) dB :
\[
S_s[k]-T[k]=3.5\ \text{dB}>0\Rightarrow D[k]=1.
\]

Si un plancher adaptatif vaut \(N_f=-50\) dB avec \(\Delta_{offset}=8\) dB :
\[
N_f+\Delta_{offset}=-42\ \text{dB}.
\]
Le seuil effectif devient alors au moins \(-42\) dB (plus strict que \(-46.5\) dB), ce qui bloque certaines fausses alarmes en bruit montant.

Enfin, exemple de mise \`a jour adaptative :
\[
N_f^{(m)}=-55,\quad N_{inst}^{(m)}=-49,\quad \alpha_{\uparrow}=0.2
\]
\[
N_f^{(m+1)}=-55+0.2\times 6 = -53.8\ \text{dB}.
\]

\section{Extraction de la cible et estimation cin\'ematique}

\subsection{Du masque binaire au pic cible}

Parmi les bins avec \(D[k]=1\), le code retient :
\[
k^\star = \arg\max_{k: D[k]=1} S_s[k].
\]
Puis il applique un filtre m\'edian sur l'historique des indices :
\[
\tilde{k}^{(m)}=\mathrm{median}\!\left(k^\star_{m-L+1},\dots,k^\star_m\right),
\]
(\(L=\)\texttt{peak\_median\_len}).

Ce filtrage est tr\`es efficace contre les sauts impulsionnels d'un bin \`a l'autre.

\subsection{Conversion fr\'equence \(\rightarrow\) vitesse}

Une fois \(\tilde{k}\) obtenu :
\[
f_{peak}=\tilde{k}\frac{F_s}{N_{FFT}},
\qquad
v_{raw}=\frac{f_{peak}}{DOPPLER\_SCALE}.
\]
En km/h :
\[
v_{raw,km/h}=3.6\,v_{raw}.
\]

\subsection{SNR de la cible}

Le SNR affich\'e est d\'efini (en dB) par :
\[
SNR_{peak}=S_s[\tilde{k}]-T[\tilde{k}].
\]
Ce n'est pas un SNR ``thermodynamique absolu'', mais un indicateur de \emph{distance au seuil adaptatif} : excellent pour le monitoring op\'erationnel.

\subsection{Filtre de Kalman : d\'erivation des \'equations utilis\'ees}

Le vecteur d'\'etat est :
\[
\mathbf{x}_m=
\begin{bmatrix}
v_m\\a_m
\end{bmatrix},
\quad
\mathbf{F}=
\begin{bmatrix}
1&dt\\0&1
\end{bmatrix},
\quad
\mathbf{H}=
\begin{bmatrix}
1&0
\end{bmatrix}.
\]

Le mod\`ele de mesure est :
\[
z_m = \mathbf{H}\mathbf{x}_m + \nu_m,\qquad \nu_m\sim\mathcal{N}(0,R).
\]

La pr\'ediction :
\[
\hat{\mathbf{x}}^-_m = \mathbf{F}\hat{\mathbf{x}}_{m-1},
\qquad
\mathbf{P}^-_m = \mathbf{F}\mathbf{P}_{m-1}\mathbf{F}^T+\mathbf{Q}.
\]

Le code impl\'emente \(\mathbf{Q}\) du mod\`ele ``acc\'el\'eration blanche'' :
\[
\mathbf{Q}=q
\begin{bmatrix}
dt^4/4 & dt^3/2\\
dt^3/2 & dt^2
\end{bmatrix}.
\]

La correction :
\[
\mathbf{K}_m=\mathbf{P}^-_m\mathbf{H}^T(\mathbf{H}\mathbf{P}^-_m\mathbf{H}^T+R)^{-1},
\]
\[
\hat{\mathbf{x}}_m=\hat{\mathbf{x}}^-_m+\mathbf{K}_m(z_m-\mathbf{H}\hat{\mathbf{x}}^-_m),
\]
\[
\mathbf{P}_m=(\mathbf{I}-\mathbf{K}_m\mathbf{H})\mathbf{P}^-_m.
\]

\paragraph{Signification des param\`etres de r\'eglage.}
\begin{itemize}[leftmargin=2em]
  \item \(q\) (\texttt{kalman\_q\_accel}) : confiance dans le mod\`ele dynamique (grands changements possibles).
  \item \(R\) (\texttt{kalman\_r\_meas}) : bruit suppos\'e de la mesure Doppler FFT.
\end{itemize}

\paragraph{Gestion des pertes de d\'etection.}
Quand aucune mesure n'est disponible :
\begin{itemize}[leftmargin=2em]
  \item pas d'update, seulement pr\'ediction,
  \item d\'ecroissance \(v\leftarrow \lambda v\) (\texttt{kalman\_miss\_decay}),
  \item annulation sous seuil (\texttt{speed\_floor}),
  \item reset complet apr\`es \texttt{kalman\_miss\_reset\_count}.
\end{itemize}
Ce bloc transforme un estimateur purement probabiliste en estimateur robuste op\'erationnel, adapt\'e \`a un capteur r\'eel sujet aux extinctions temporaires.

\subsection{Exemple num\'erique (section 8)}

\paragraph{1) Fr\'equence et vitesse.}
Avec \(N_{FFT}=8192\), \(\Delta f=12.207\) Hz/bin.  
Supposons \(\tilde{k}=132\) :
\[
f_{peak}=132\times 12.207 = 1611.3\ \text{Hz}.
\]
\[
v_{raw}=\frac{1611.3}{161}=10.01\ \text{m/s}=36.0\ \text{km/h}.
\]

\paragraph{2) SNR au pic.}
Si \(S_s[\tilde{k}]=-37.2\) dB et \(T[\tilde{k}]=-49.0\) dB :
\[
SNR_{peak}=11.8\ \text{dB}.
\]

\paragraph{3) Pas Kalman (illustratif).}
Prenons \(dt=0.01024\) s, estimation pr\'edite \(v^- = 9.6\) m/s, mesure \(z=10.0\) m/s.
Si le gain vitesse \(K_0=0.35\), alors :
\[
v = v^- + K_0(z-v^-)
= 9.6 + 0.35(0.4)
= 9.74\ \text{m/s}.
\]
Le filtre corrige partiellement vers la mesure, sans la suivre brutalement.

\section{Publication des r\'esultats et rendu GUI}

\subsection{Donn\'ees partag\'ees}

Le thread DSP publie (mutex) :
\begin{itemize}[leftmargin=2em]
  \item ADC glissant brut,
  \item spectre dB, seuil CFAR, masque d\'etections,
  \item \(f_{peak}\), SNR, vitesse brute et Kalman,
  \item plancher de bruit adaptatif,
  \item spectrogramme circulaire,
  \item FPS, \'etat connexion, pertes de trames.
\end{itemize}

\subsection{Sorties visuelles}

Le GUI construit :
\begin{itemize}[leftmargin=2em]
  \item \textbf{Oscilloscope} (temps), avec fen\^etre temporelle ajustable,
  \item \textbf{FFT + CFAR} (fr\'equence), avec curseur et marqueur pic,
  \item \textbf{Spectrogramme waterfall} (temps-fr\'equence),
  \item \textbf{Panneau num\'erique} (vitesse, SNR, statut cible),
  \item \textbf{Trend Doppler} long terme.
\end{itemize}

Le spectrogramme est stock\'e en tampon circulaire 2D, texture OpenGL mise \`a jour en filtrage \texttt{NEAREST} pour une lecture nette des bins.

\section{Lecture ``professeur'' : compromis et limites}

\subsection{R\'esolution et latence}

La r\'esolution fr\'equentielle vaut :
\[
\Delta f = \frac{F_s}{N_{FFT}}
\]
et la r\'esolution vitesse :
\[
\Delta v = \frac{\Delta f}{161}
\]

Augmenter \(N_{FFT}\) am\'eliore \(\Delta f\), mais augmente co\^ut de calcul et inertie d'estimation. Le hop \`a 10.24 ms maintient une cadence de rafra\^ichissement \'elev\'ee.

\subsection{Ce que CFAR r\'esout, et ce qu'il ne r\'esout pas}

CFAR traite bien l'h\'et\'erog\'en\'eit\'e du bruit local, mais :
\begin{itemize}[leftmargin=2em]
  \item il reste sensible \`a des interférences structur\'ees (secteur, parasites harmoniques),
  \item il peut d\'egrader en multi-cibles tr\`es proches si les fen\^etres de garde/r\'ef\'erence sont mal r\'egl\'ees.
\end{itemize}

D'o\`u l'int\'er\^et des briques compl\'ementaires (notch, adaptive floor, median peak, Kalman).

\subsection{Pourquoi vitesse brute + vitesse filtr\'ee}

La vitesse brute montre le comportement instantan\'e physique; la vitesse Kalman fournit une estimation op\'erationnelle plus stable (pilotage UI, d\'ecision humaine).

\section{Conclusion}

\texttt{pc\_analyzer} r\'ealise une cha\^ine DSP radar compl\`ete et moderne :
\begin{enumerate}[leftmargin=2em]
  \item acquisition distribu\'ee (STM32 + UDP fragment\'e),
  \item reconstruction glissante et pr\'e-traitements robustes,
  \item analyse spectrale haute r\'esolution,
  \item d\'etection adaptative (CFAR + bruit adaptatif + mains/notch),
  \item estimation cin\'ematique stabilis\'ee (Kalman),
  \item restitution visuelle et m\'etrologique temps r\'eel.
\end{enumerate}

Autrement dit, le logiciel n'est pas une simple FFT affich\'ee : c'est un estimateur radar complet, con\c{c}u pour transformer un flux ADC bruit\'e en information vitesse exploitable.

\end{document}

